
\documentclass[landscape]{article}
\usepackage[a4paper,margin=1cm]{geometry}
\usepackage{multicol}
\usepackage{pdflscape}
\usepackage{amsmath,amssymb}
\usepackage{tcolorbox}
\tcbuselibrary{skins,breakable}
\setlength{\columnsep}{0.8cm}

\begin{document}
\begin{landscape}
\begin{multicols}{4}

Ya estamos en condiciones de definir uno de los conceptos más importantes asociados a los polinomios: el de cero o raíz de un polinomio. El cálculo de los ceros de un polinomio será de suma utilidad para resolver ecuaciones, inecuaciones y cualquier tipo de problema que requiera la factorización de polinomios.

\begin{tcolorbox}[colback=blue!5!white,colframe=blue!75!black,title=\textbf{Definición (Cero o raíz de un polinomio)}]
Decimos que un número real $\alpha$ es un cero o raíz del polinomio $p(x)$ si $p(\alpha) = 0$.
\end{tcolorbox}

\noindent\textbf{Ejemplo 1.10.} Calcule el valor numérico del polinomio $p(x) = -x^3 + 4x - 5$ para $x = 3$.

\noindent\textbf{Solución:}
$p(3) = -(3)^3 + 4(3) - 5 = -27 + 12 - 5 = -20$.

\columnbreak

\noindent\textbf{Ejemplo 1.11.} Verifique que $-2$ es una raíz del polinomio $p(x) = x^3 + x^2 - 5x - 6$.

\noindent\textbf{Solución:}
$p(-2) = (-2)^3 + (-2)^2 - 5(-2) - 6 = -8 + 4 + 10 - 6 = 0$,
por lo que $p(-2)=0$ confirma que $x=-2$ es una raíz.

\noindent\textbf{Ejemplo 1.12.} Determine el valor de $\alpha$ para que $4$ sea raíz del polinomio $p(x) = -x^3 + \alpha x^2 + 2x + 6$.

\noindent\textbf{Solución:} Para que $4$ sea una raíz de $p(x)$, debe ocurrir $p(4) = 0$. Luego:

\begin{align*}
0 &= -(4)^3 + \alpha (4)^2 + 2(4) + 6,\\
0 &= -64 + 16\alpha + 8 + 6,\\
0 &= -50 + 16\alpha,\\
16\alpha &= 50,\\
\alpha &= \frac{50}{16} = \frac{25}{8}.
\end{align*}

\columnbreak

Recordemos que un polinomio $p(x)$ es divisible por otro polinomio $q(x)$ si el resto de dividir $p$ por $q$ es igual a cero. El siguiente teorema muestra que cuando el divisor es de la forma $q(x) = x - k$, podemos determinar si $p(x)$ es divisible por $q(x)$ sin necesidad de realizar la división: basta con calcular $p(k)$.

\begin{tcolorbox}[colback=red!5!white,colframe=red!75!black,title=\textbf{Teorema del resto}]
El resto de dividir $p(x)$ por un polinomio de la forma $x - k$ (con $k$ un número real cualquiera) es igual a $p(k)$.

\medskip\noindent\textbf{Demostración:} Si $p(x)$ es un polinomio de grado mayor o igual que 1, por el Algoritmo de la División existen polinomios $c(x)$ y $r(x)$ tales que:
\begin{equation*}
p(x) = (x - k)c(x) + r(x),
\end{equation*}
donde $r(x)$ es el resto y cumple $\deg r < \deg(x - k) = 1$. Evaluando en $x = k$ se tiene:
\begin{equation*}
p(k) = (k - k)c(k) + r = r.
\end{equation*}
\end{tcolorbox}

\noindent\textbf{Ejemplo 1.13.} Calcule el resto de dividir $p(x) = -x^4 - 2x^3 + 3x^2 + x + 1$ por $q(x) = x + 1$.

\noindent\textbf{Solución:} 
$p(-1) = -(-1)^4 - 2(-1)^3 + 3(-1)^2 + (-1) + 1 = -1 + 2 + 3 - 1 + 1 = 4$.
Por lo tanto, el resto es $4$.

\columnbreak

\noindent\textbf{Ejemplo 1.14.} Halle el valor de $a \in \mathbb{R}$ para que, al dividir $p(x) = x^4 + 3x^3 + ax^2 - 13$ por $x + 3$, el resto sea igual a $5$.

\noindent\textbf{Solución:}
\begin{align*}
(-3)^4 + 3(-3)^3 + a(-3)^2 - 13 &= 5\\
81 - 81 + 9a - 13 &= 5\\
9a &= 18\\
a &= 2.
\end{align*}

Del Teorema del Resto se deduce el siguiente resultado:

\begin{tcolorbox}[colback=red!5!white,colframe=red!75!black,title=\textbf{Teorema}]
Si $\alpha \in \mathbb{R}$ es una raíz del polinomio $p(x)$, entonces $p(x)$ es divisible por $x - \alpha$.
\end{tcolorbox}

\noindent\textbf{Ejercicio 1.11.} Determine cuáles de los siguientes polinomios son divisibles por $x + 2$:

\begin{enumerate}
\item $p(x) = -x^3 - 5x + 2$,
\item $p(x) = -x^3 + 5x + 2$,
\item $p(x) = 2x^4 + x^3 - 3$,
\item $p(x) = x^4 + x^2 + 12x + 4$.
\end{enumerate}

\begin{tcolorbox}[colback=blue!5!white,colframe=blue!75!black,title=\textbf{Definición (Polinomio irreducible)}]
Un polinomio $p(x)$ de grado mayor o igual que 1 se dice irreducible si cualquier factorización 
$p(x) = m(x)n(x)$ implica que $\deg m = 0$ o $\deg n = 0$.
\end{tcolorbox}

Un polinomio con coeficientes reales es irreducible en $\mathbb{R}$ si y sólo si es de grado 1 o es de grado 2 sin raíces reales.

\begin{observacion}[1.6] De aquí en adelante, \emph{irreducible} significará irreducible en $\mathbb{R}$. \end{observacion}

\noindent\textbf{Ejemplo 1.15.} $p(x) = 3x - 5$ y $q(x) = x^2 + 1$ son irreducibles.

\noindent\textbf{Ejemplo 1.16.} $p(x) = x^3 + x$ no es irreducible pues $p(x)=x(x^2+1)$.

\noindent\textbf{Ejemplo 1.17.} $p(x) = x^2 - 3x - 4 = (x-4)(x+1)$ no es irreducible.

\begin{tcolorbox}[colback=blue!5!white,colframe=blue!75!black,title=\textbf{Definición (Factorización de un polinomio)}]
Factorizar $p(x)$ es escribirlo como producto de polinomios irreducibles.
\end{tcolorbox}

\noindent\textbf{Ejemplo 1.18.} 
$p(x) = 2x^3 + 3x^2 - 2x - 3 = 2(x-1)(x+1)\bigl(x+\tfrac32\bigr)$.

\begin{observacion}[1.7] Las raíces son $1,-1,-\tfrac32$. \end{observacion}

\end{multicols}
\end{landscape}
\end{document}
